\section{Using the engine from Python}
\label{sec:library}

The engine is a library first and a program second.  Everything the command line
does is available as a function, and everything it prints is available as data.

\subsection{The public surface}

\begin{longtable}{@{}p{4.6cm}p{9.8cm}@{}}
\toprule
Name & Purpose \\
\midrule
\endhead
\texttt{Cube} & an immutable cube state: \texttt{apply}, \texttt{apply\_moves},
\texttt{apply\_sequence}, \texttt{apply\_scramble}, \texttt{is\_solved},
\texttt{is\_solvable}, \texttt{cubie}, \texttt{to\_dict} \\
\texttt{CubieCube} & the same cube at the piece level: corner and edge permutations
with orientations, plus \texttt{validate} \\
\texttt{MOVES}, \texttt{get\_move}, \texttt{invert} & the 54 canonical moves and
their inverses \\
\texttt{parse\_moves}, \texttt{parse\_scramble}, \texttt{format\_moves},
\texttt{normalize\_sequence} & notation in and out \\
\texttt{generate\_scramble}, \texttt{generate\_scrambles} & scrambles, with a seed
and a style \\
\texttt{Scramble} & the movement list, its text, the state it produces, its seed \\
\texttt{solve\_cube} & the CFOP pipeline: the whole solution as a \texttt{Solution} \\
\texttt{cross\_options} & the six cross colours with their lengths, ranked \\
\texttt{Solution}, \texttt{Stage}, \texttt{CrossOption} & the result objects, each
with \texttt{to\_dict} \\
\texttt{solve}, \texttt{SearchLimits} & the two-phase solver on its own, with depth
and time limits \\
\texttt{CancelToken} & cooperative cancellation for a solve \\
\texttt{preload} & build the scrambler tables now \\
\texttt{net\_layout}, \texttt{render\_ansi}, \texttt{render\_emoji},
\texttt{render\_letters}, \texttt{NetLayout}, \texttt{DEFAULT\_PALETTE} & the
unfolded net, as data or as text \\
\texttt{RubiksError} and subclasses & structured errors, each with
\texttt{to\_dict} \\
\bottomrule
\end{longtable}

\subsection{Scramble and solve}

\begin{lstlisting}[style=python]
import rubiks_cfop as cube

scramble = cube.generate_scramble(seed=42)
print(scramble.text)                 # "D B2 F2 ..."
print(scramble.style.value)          # "random-state"
print(scramble.state.is_solvable())  # True

solution = cube.solve_cube(scramble.state, scramble=scramble.text)
for stage in solution.stages:
    print(f"{stage.label}: {stage.text}")
print(solution.total_moves, solution.verified)
\end{lstlisting}

\begin{lstlisting}[style=shell]
Cross: z z R' F D' F2 R' F
F2L1: U2 B' R2 B R2
F2L2: F' L F L2 U' L
F2L3: U L U' B U B' L'
F2L4: U B U B' U2 R' U2 R
OLL (40): U R' F R U R' U' F' U R
PLL (F Perm): U' L R2 U R U R2 U' R' U' R2 U' R U2 L' U R' U2
62 True
\end{lstlisting}

\subsection{Choosing the cross}

\begin{lstlisting}[style=python]
solution = cube.solve_cube(cube.Cube.solved().apply_scramble("R U R' U'"), cross="white")

for option in cube.cross_options(scramble.state):
    print(option.color.value, option.length)
\end{lstlisting}

The colour may be a \texttt{Color} or its name; \texttt{"best"} is the default.

\subsection{Drawing the cube}

\begin{lstlisting}[style=python]
layout = cube.net_layout(scramble.state)   # 9 rows x 12 columns of Color
print(cube.render_emoji(layout))           # for a terminal without colour
print(cube.render_ansi(layout))            # 24-bit ANSI background colours

for face, rows in layout.to_dict()["faces"].items():
    print(face, rows[0])                   # the top row of each face, by name
\end{lstlisting}

\texttt{NetLayout} holds colours, not escape codes: an interface draws the cube in
its own way from the same structure.  \texttt{render\_letters} exists for logs and
diffs, and is deliberately not the default, because the program's output is meant to
show real colours.

\subsection{Cancelling a solve}

A solve can take a few seconds, and an interactive program must be able to abandon
one.

\begin{lstlisting}[style=python]
import threading
import rubiks_cfop as cube

cancel = cube.CancelToken()

def work():
    try:
        solution = cube.solve_cube(scramble.state)
    except cube.CancelledError:
        print("abandoned")

thread = threading.Thread(target=work)
thread.start()
# ... later, from the user interface thread:
cancel.cancel()
\end{lstlisting}

Passing a token to a solve is a matter of scheduling: the pipeline checks it at
regular intervals and raises \texttt{CancelledError} promptly.  The two-phase solver
accepts the same token and an optional timeout through \texttt{SearchLimits}.

\subsection{Errors}

Every error derives from \texttt{RubiksError} and carries a stable \texttt{code}, a
message and structured details, so a caller never has to parse a string:

\begin{center}
\begin{tabular}{@{}ll@{} }
\toprule
Class & When \\
\midrule
\texttt{NotationError} & a move sequence that cannot be parsed, with
\texttt{token} and \texttt{index} \\
\texttt{InvalidMoveError} & a single move that is not valid on a $3\times3\times3$ \\
\texttt{InvalidStateError} & a malformed state: wrong length, wrong counts, moved
centres, a piece that cannot exist \\
\texttt{UnsolvableStateError} & a well formed but impossible state: a twisted
corner, a flipped edge, wrong parity \\
\texttt{SearchLimitError} & a search ran out of depth budget \\
\texttt{SolveTimeoutError} & a solve ran out of time \\
\texttt{CancelledError} & a solve was cancelled \\
\bottomrule
\end{tabular}
\end{center}

\begin{lstlisting}[style=python]
try:
    cube.Cube.solved().apply_sequence("R Q")
except cube.NotationError as error:
    print(error.to_dict())
    # {'code': 'notation-error', 'message': 'not WCA move notation',
    #  'details': {'index': 1, 'token': 'Q'}}
\end{lstlisting}

\subsection{Threading and immutability}

\texttt{Cube}, \texttt{CubieCube}, \texttt{Move}, \texttt{Solution} and
\texttt{Stage} are frozen value objects: sharing one between threads is safe, and
applying a move returns a new cube rather than changing one.  Solved data is never
mutated, the caches are read-only tables, and the solvers keep their state in local
variables, so a program may solve on a worker thread while its interface keeps
running.  The only expensive one-off work --- the scrambler tables --- is exposed as
\texttt{preload()} precisely so that it can be moved off the interactive path.
